\scanpage{54}
\[
(31)\quad\left\{\begin{aligned}
\varphi_{R_i^\omega,\omega i}&:\Pi_{\omega i}\xrightarrow{\sim}\Pi_1(\widetilde\Sigma,R_i^\omega)\\
&\text{or }\varphi_{\omega i}\\
&\text{or }c_{\omega i}^{\omega'}\\
\varphi_{S_i^\varepsilon,\omega\deleted{\unsure{i}}}&:\Pi_\omega\xrightarrow{\sim}\Pi_1(\widetilde\Sigma,S_i^\varepsilon)\\
&\text{or }\varphi_\omega\\
&\text{or }d_i^{\varepsilon}
\end{aligned}\right.
\]
It seems to me most convenient to write simply $\varphi_i,\varphi_\varepsilon$ to specify the \unsure{standard|standard} $\Pi_1$ group $\Pi_i$ or $\Pi_\varepsilon$ which are sources, the base point $x$ ($=R_i^\varepsilon$ or $S_i^\varepsilon$, say) not being in doubt in the calculation.
\deleted{The arcs $a_i^\omega$ are:}
\deleted{The arcs $a_i^\omega$ are visibly too long, they must be cut in two by setting
\[
(31)\quad\left\{\begin{aligned}
a_i^\omega&=a_i^{\omega-}a_i^{\omega+}\\
a_i^{\omega+}&:R_i^\omega\longrightarrow S_i^\omega\qquad a_i^{\omega-}:S_i^\omega\longrightarrow R_i^{\omega'}
\end{aligned}\right.
\]
}
Finally, for formal calculations, one must give a way of identifying the $\Pi_i,\Pi_\omega$ with a $=$ \unsure{standard|standard} group \deleted{$\Pi_{0,3}$}, for which we shall take
\[
\begin{aligned}
(32)\qquad\Pi_{0,3}&=\Pi_1\left(\deleted{\Sigma}\overbrace{\mathbb P^1_{\mathbb C}\setminus\{0,1,\infty\}}^{U_{0,3}(\mathbb C)},P^+\right)\\
&=\{l_0,l_1,l_\infty\mid l_\infty l_1l_0=1\}\\
&\subset\widetilde\Pi_{03}=\Pi_1(U_{03}(\mathbb C),\mathcal G_3^*\inserted{\times\{1,\tau\}},P^+)
\end{aligned}
\]
\[
=\{\rho,\sigma_\infty,\widetilde\sigma_\infty\mid \rho^3=\sigma_\infty^2\inserted{=\widetilde\sigma_\infty^2}=1,\quad \widetilde\sigma_\infty(\rho)=\rho^{-1},\quad\widetilde\sigma_\infty(\widetilde\sigma_\infty)=1\}
\]
\inserted{i.e. $(\sigma_\infty\widetilde\sigma_\infty)^2=1$}
\deleted{$=\{\rho,\sigma_\infty\mid\rho^3=\sigma_\infty^2=1\}\simeq\operatorname{SL}(2,\mathbb Z)/\{\pm1\}$}
\inserted{$\simeq\operatorname{GL}(2,\mathbb Z)/\{\pm1\}$}
\margin{53}
